\section{How much can the height change when the word changes?}

% BEGIN H89V1
\paragraph{H89V1. Cancellation of the level and the common normalizations.}
Fix a member of a parameter family at scale $x>x_1$. A word $u$
chooses a period map, with assigned logarithmic height allowance $h_u$
and level prime $T_u$. Write $A=\sum_{s\in TT}a_s$ and
$H_u=h_u+A\log T_u$. The assertion needed in the ninth height
estimate is a bound for the diameter of the set of numbers $H_u$.
The common origin from which the heights are measured can be
discarded exactly.

In fact let $m$ be the number of lifts and $f$ the number of
archimedean labels $\tau$. Write $\lambda$ for the member's
parameter $\log\Omega$, and put
\[
  g(y)=y\log y-y-\lambda y .
\]
With $N_i=k+i(t+1)$ and $y_{u,i,\tau,r}$ the row exponent at
the $(i+1)$st lift, there is a real number $H_*$ for which
\[
 H_u=H_*+\sum_{i=0}^{m-1}\sum_\tau
                \sum_{r=0}^{N_i-1}g(y_{u,i,\tau,r}).
\]
To see this, use the defining height formula of the given family.
Its level term is $-\sum_s a_s\log T_u$. Its archimedean kernel
term at each lift and place is the sum of $y\log y-y$ minus
$\lambda\sum y$. The remaining terms are the common size allowances,
the lower hook-product terms, and the mean logarithm of the metric
normalization. None involves $u$. This proves the displayed
identity for the actual $h_u$. It makes no assertion about the
absolute height of an individual map.
% END H89V1

% BEGIN H89V2
\paragraph{H89V2. Only the interior of the first weight block moves.}
Suppress $\tau$ for a moment. The first $k$ row exponents have
the form
\[
 y_{u,i,r}=\alpha_{u,r}+e_i-N_i-t/2+r,\qquad 0\le r<k .
\]
The two ends $\alpha_{u,0},\alpha_{u,k-1}$ are prescribed
independently of $u$. The remaining rows come in blocks indexed
by $0\le q<i$, with the common value $e_i-e_q-(i-q)(t+1)$
throughout a block. Thus in the entire sum only the first block's
interior can change. Every moving coordinate satisfies
$0\le\alpha_{u,r}\le S$, where $S$ is the specified width of
the weight interval. In particular its displacement between two
words is at most $S$.

For completeness, each moving $y$ is at least one. The
$\alpha$'s are strictly decreasing half-integers in the same
coset of $\mathbb Z$, so $\alpha_{u,r}+r$ is nonincreasing in
$r$. The gap conditions give $e_{i+1}-e_i\ge t+3$, while
$N_{i+1}-N_i=t+1$. Hence the minimum head value occurs at
$i=0,r=k-1$ for the word and place in question. This value is
$\alpha_{u,k-1}+e_0-1-t/2>0$, since $\alpha_{u,k-1}>-b_0$
and $e_0-b_0\ge(t+3)/2$. Here $b_0$ is the first weight of the
interlacing parameter. Integrality follows from odd $t$, integral
$e_i$ and half-integral $\alpha_{u,r}$. The constant tail blocks
are positive as well by the same gap between consecutive $e_i$.

If $S+e_i\le M$ for each lift and place, $M\ge1$, the
moving exponents lie in $[1,M]$, since $N_i+t/2-r\ge0$.
On this interval $|g'(y)|=|\log y-\lambda|
\le\log M+|\lambda|$. The mean value theorem in each moving
coordinate now gives, simultaneously for all words,
\[
            |H_u-H_{u'}|\le mfkS\,(\log M+|\lambda|).
\]
The estimate costs at most $k$ variations per lift and place
even when $N_i$ is much larger than $k$.
% END H89V2

% BEGIN H89V3
\paragraph{H89V3. Uniform diameter bound on the original families.}
Take a family in the selected class of P88A1 and put $\ell=x^\epsilon$, $b=e^\ell$, $X'=\exp(x^{\epsilon/3})$
and $W=b/(nX')$. We claim that for every sufficiently
late member, with $B_0=fmb\ell$,
\[
       |H_u-H_{u'}|\ \le\ \frac{4B_0}{nX'} .
\]
Here $u,u'$ are any two words of the member. The width and spacing assumptions
give $kS\le W$ and $t\eta\le W/x$, where $\eta$ is the width
of the lift-weight windows. In particular $S,\eta,W\le b$.

For the original three weight windows H88E1 gives $e_i\le3b$
at every place and lift. Together with $S\le b$ this gives the
single upper endpoint $M=4b$ in H89V2.

The normalization is
$\lambda=\gamma+\log(2\pi)-\log\operatorname{rd}(E)-d_1$,
where $d_1$ is the prescribed logarithmic derivative at one.
The size assumptions give $|d_1|\le\ell$ and
$1\le\operatorname{rd}(E)\le b$. For the latter upper bound,
the root discriminant is at most the conductor parameter
$Q_S qd$, whose $36$th power times $t^{150}$ is at most
$\sqrt b$. Consequently $|\lambda|\le 2\ell+6$.
For fixed $\epsilon>0$ choose a threshold so that
$\log 4+6\le x^\epsilon$ for every larger $x$. This gives
$\log M+|\lambda|\le4\ell$. H89V2 and $kS\le W$ give the
claimed bound.

The threshold can be increased to exceed the family's $x_1$.
The family identifies its scale and exponent with $x,\epsilon$
and supplies the height formula and size inequalities for every
index there. Since $nX'\ge1$, $m\ge1$ and
$a\Lambda_a\ge0$ (here $a=xb$ and
$\Lambda_a=\log(a\operatorname{rd}(E))$), the estimate implies
\[
 |(h_u+A\log T_u)-(h_{u'}+A\log T_{u'})|
          \le 4\bigl(|TT|a\Lambda_a+mB_0\bigr).
\]
This is the full ninth entry with one constant and one cutoff
before the choice of the index or either word. On intersecting
with the range of P88A1 one can enlarge the common constant,
since $B_0$ and $|TT|a\Lambda_a+mB_0$ are nonnegative.
% END H89V3
